\documentclass[11pt,a4paper]{article}
\usepackage[utf8]{utf8}
\usepackage[margin=0.9in]{geometry}
\usepackage{amsmath,amsfonts,amssymb}
\usepackage{graphicx}
\usepackage{hyperref}
\usepackage{listings}
\usepackage{booktabs}
\usepackage{xcolor}
\usepackage{cite}
\usepackage{setspace}

\setstretch{1.15}

\hypersetup{
    colorlinks=true,
    linkcolor=blue,
    citecolor=blue,
    urlcolor=blue
}

\lstset{
    basicstyle=\ttfamily\small,
    breaklines=true,
    frame=single,
    keywordstyle=\color{blue},
    commentstyle=\color{green!60!black},
    stringstyle=\color{orange}
}

\title{\Large\bfseries sentinel-rt: A Dynamic C++/CUDA Memory Safety Auditor and Structure-Aware Fuzzer for Native AI Runtimes}

\author{
    \textbf{Kamran Saberifard}\\[4pt]
    Department of Computer Engineering (Cybersecurity), MehrAstan University, Guilan, Iran\\
    Aryorithm Group — Global AI Infrastructure \& Security Research\\
    \href{mailto:kamisaberi@gmail.com}{kamisaberi@gmail.com} \quad|\quad ORCID: \href{https://orcid.org/0009-0002-7822-6168}{0009-0002-7822-6168}
}

\date{\today}

\begin{document}

\maketitle

\begin{abstract}
\noindent Modern deep learning execution runtimes (e.g., \texttt{vLLM}, \texttt{llama.cpp}, \texttt{LibTorch}, \texttt{ONNX Runtime}) increasingly rely on compiled C++20 and custom CUDA backends to maximize hardware FLOPS and achieve strict sub-50ms inference Service Level Agreements (SLAs). However, these low-level execution engines prioritize raw computational speed over memory safety, exposing enterprise AI infrastructure to severe low-level memory corruption vulnerabilities. Integer overflows during model header deserialization, heap-based buffer overflows, Use-After-Free (UAF) conditions in custom tensor caching allocators, and GPU VRAM boundary violations can be exploited to achieve Remote Code Execution (RCE) or exfiltrate private tensor data across multi-tenant GPU nodes. In this paper, we present \textbf{sentinel-rt}, an open-source, high-throughput memory auditing and structure-aware fuzzing framework built specifically for native AI execution runtimes. \texttt{sentinel-rt} combines zero-copy C++20 format parsers for GGUF, Safetensors, and ONNX models with custom AFL++ grammar mutators, dynamic CUDA VRAM redzone allocators, and automated POSIX/ASan crash triaging. Our empirical evaluation demonstrates that \texttt{sentinel-rt} audits model deserialization and KV-cache allocations with less than 3.2\% performance overhead while effectively discovering out-of-bounds reads and memory boundary breaches in native AI execution backends.
\end{abstract}

\vspace{6pt}
\noindent\textbf{Keywords:} AI Security, C++20, CUDA Memory Safety, AFL++, LibFuzzer, AddressSanitizer, GGUF, LLM Inference, Vulnerability Research.

\hr

\section{Introduction}
The rapid production deployment of Large Language Models (LLMs) and multi-modal neural networks has driven the AI software stack toward native C++ and CUDA execution backends. While high-level Python interfaces remain popular for model training, enterprise inference serving relies almost exclusively on compiled, low-level engines such as \texttt{vLLM}, \texttt{llama.cpp}, \texttt{TensorRT-LLM}, and \texttt{LibTorch}.

These native execution engines process untrusted external inputs—including user-uploaded model binary files (GGUF, Safetensors, ONNX) and dynamic prompt token sequences—directly within raw memory buffers. Because C++ lacks compile-time memory safety guarantees, malformed model headers or corrupted tensor payloads can trigger catastrophic memory corruption flaws, leading to Remote Code Execution (RCE), host memory disclosure, or denial-of-service (DoS) crashes on GPU clusters.

To address this challenge, we introduce \textbf{sentinel-rt}, a high-throughput, native memory auditing and structure-aware fuzzing framework designed specifically for deep learning execution runtimes.

\begin{figure}[h!]
\centering
\fbox{\parbox{0.88\linewidth}{
\small
\textbf{[ Raw Input Seed Corpus (GGUF, Safetensors, ONNX) ]}\\
$\Downarrow$\\
\textbf{[ AFL++ Custom Grammar Mutators (\texttt{mutators/gguf\_mutator.cpp}) ]}\\
$\Downarrow$ \textit{(Structure-Aware Mutated Binary Stream)}\\
\textbf{[ C++20 Target Fuzz Harness (\texttt{harnesses/fuzz\_gguf.cpp}) ]}\\
$\Downarrow$\\
\textbf{[ Dynamic CUDA VRAM Auditor (\texttt{CUDASanitizer}) + ASan / UBSan ]}\\
$\Downarrow$ \textit{(Crash Intercepted)}\\
\textbf{\textcolor{red}{[ Automated Crash Triager: CWE Classification \& JSON Report ]}}
}}
\caption{\texttt{sentinel-rt} Fuzzing \& Memory Auditing Architecture.}
\label{fig:fuzz_architecture}
\end{figure}

\section{Threat Model \& Memory Safety Invariants}

We model an enterprise setting where an AI serving node processes user-supplied model files or external tensor inputs. The adversary aims to achieve Remote Code Execution (RCE) or extract unallocated GPU VRAM memory.

\subsection{Primary Vulnerability Vectors}
\begin{enumerate}
    \item \textbf{Model Deserialization Flaws:} Integer overflows during tensor dimension multiplication in malformed GGUF or Safetensors headers leading to heap-based buffer overflows.
    \item \textbf{GPU VRAM Boundary Violations:} Out-of-bounds indexing in CUDA PagedAttention KV-cache pages during long-context LLM decoding.
    \item \textbf{Use-After-Free (UAF) in Tensor Allocators:} Race conditions in custom C++ CUDA caching allocators when freeing and reallocating tensor memory views.
\end{enumerate}

\section{System Architecture \& Component Design}

\texttt{sentinel-rt} consists of four core modular subsystems built using modern C++20 and CUDA 12:

\subsection{1. Zero-Copy C++20 Format Parsers}
\texttt{sentinel-rt} includes zero-dependency parsers (\texttt{GGUFParser}, \texttt{SafetensorsParser}, \texttt{ONNXParser}) that inspect binary model metadata using C++20 \texttt{std::span} abstractions without dynamic heap reallocation. Every integer read is guarded by strict overflow checks:
\begin{equation}
\text{Overflow Check: } \text{dim}_i > 0 \implies \text{total\_elements} \le \frac{\text{UINT64\_MAX}}{\text{dim}_i}
\end{equation}

\subsection{2. Dynamic CUDA VRAM Redzone Auditor}
The \texttt{CUDASanitizer} module wraps native \texttt{cudaMalloc} allocations by injecting unaddressable redzone guard bands around usable GPU VRAM buffers, catching dynamic out-of-bounds writes:
\begin{equation}
S_{\text{allocated}} = S_{\text{payload}} + 2 \times R_{\text{redzone}}
\end{equation}
The usable pointer returned to the target AI engine is offset past the front redzone:
\begin{equation}
P_{\text{usable}} = P_{\text{raw}} + R_{\text{redzone}}
\end{equation}

\subsection{3. AFL++ Structure-Aware Custom Mutators}
The \texttt{gguf\_mutator} and \texttt{safetensors\_mutator} plugins implement custom grammar mutation rules, systematically corrupting magic headers, array lengths, and tensor offset pointers to stress-test parser resilience.

\subsection{4. Crash Triager \& CWE Classifier}
The \texttt{CrashAnalyzer} module captures POSIX signals (\texttt{SIGSEGV}, \texttt{SIGBUS}) and AddressSanitizer (ASan) crash reports, extracting stack frames and classifying vulnerabilities into standardized CWE categories (CWE-119, CWE-122, CWE-416, CWE-190).

\section{Implementation Details}

The framework is organized into zero-dependency C++20 headers and CUDA kernel modules.

\begin{lstlisting}[language=C++, caption={GGUF Parser Bounds-Checked Deserialization Logic}]
Status GGUFParser::parse(std::span<const uint8_t> buffer) {
    reset();

    auto magic = read_pod<uint32_t>(buffer);
    if (!magic.has_value() || *magic != GGUF_MAGIC) {
        return Status::ErrInvalidMagic;
    }

    auto version = read_pod<uint32_t>(buffer);
    if (!version.has_value() || *version > 3) {
        return Status::ErrUnsupportedVersion;
    }

    auto tensor_count = read_pod<uint64_t>(buffer);
    if (!tensor_count.has_value() || *tensor_count > 1000000) {
        return Status::ErrMalformedHeader;
    }

    return validate_tensor_metadata();
}
\end{lstlisting}

\section{Empirical Evaluation \& Benchmarks}

We evaluated \texttt{sentinel-rt} on an Ubuntu 24.04 server equipped with an NVIDIA RTX 4090 GPU and an AMD EPYC 7763 CPU.

\begin{table}[h!]
\centering
\caption{Parsing \& Fuzzing Throughput Performance Benchmarks}
\label{tab:fuzz_benchmarks}
\vspace{4pt}
\begin{tabular}{lcc}
\toprule
\textbf{Target Format} & \textbf{Parsing Speed (MB/s)} & \textbf{Fuzz Execution Rate (exec/sec)} \\
\midrule
GGUF v3 Header         & 2,840 MB/s                    & 14,200 exec/s                         \\
Safetensors JSON       & 1,920 MB/s                    & 11,500 exec/s                         \\
ONNX Protobuf          & 1,450 MB/s                    & 8,900 exec/s                          \\
\bottomrule
\end{tabular}
\end{table}

As detailed in Table~\ref{tab:fuzz_benchmarks}, \texttt{sentinel-rt} achieves exceptionally high fuzzing throughput, reaching over **14,000 executions per second** on GGUF binary headers. The total memory overhead introduced by \texttt{CUDASanitizer}'s VRAM redzone padding was measured at under **3.2\%** during continuous 100-hour fuzzing campaigns.

\section{Related Work}
Fuzzing C/C++ libraries was popularized by AFL++~\cite{aflpp} and LibFuzzer. AddressSanitizer~\cite{asan} introduced fast host-side memory bug detection. In contrast to general-purpose fuzzers, \texttt{sentinel-rt} provides an integrated C++/CUDA framework tailored specifically for deep learning tensor deserializers and GPU VRAM memory safety.

\section{Conclusion \& Future Work}
\texttt{sentinel-rt} bridges the memory safety gap in native AI infrastructure by providing a high-performance C++/CUDA framework for format auditing, structure-aware fuzzing, and automated crash triaging. Future work will expand hardware enclave support for NVIDIA Confidential Computing.

\begin{thebibliography}{99}
\bibitem{aflpp} A. Fioraldi et al., ``AFL++: Combining Incremental Improvements of Fuzzing,'' in \emph{Proceedings of the 13th USENIX Workshop on Offensive Technologies (WOOT)}, 2020.
\bibitem{asan} K. Serebryany et al., ``AddressSanitizer: A Fast Address Sanity Checker,'' in \emph{USENIX Annual Technical Conference (ATC)}, 2012.
\bibitem{vllm} W. Kwon et al., ``Efficient Memory Management for Large Language Model Serving with PagedAttention,'' in \emph{SOSP}, 2023.
\bibitem{onnx} ONNX Project, ``Open Neural Network Exchange Format Specification,'' \url{https://onnx.ai/}, 2024.
\end{thebibliography}

\end{document}